% Originally: content/week-04.tex, \section{Chain rule} (Corsin Week 4, Monday)

% Source: Corsin Nick/Class Notes/Week 4.pdf, p. 5
\section{Chain rule}
\label{sec:chain_rule}

\begin{theorem}[Chain rule]
\label{thm:chain_rule}
The \newterm{chain rule} \germanfor{Kettenregel} states: for
$f : U \subseteq \mathbb{R}^n \to \mathbb{R}^m$, $g : V \subseteq \mathbb{R}^m \to \mathbb{R}^d$
differentiable, it holds:
\[D(g\circ f)(x) = Dg(f(x))\,Df(x).\]
The same identity holds for the corresponding Jacobi matrices, where the composition on the
right becomes an ordinary matrix product.
\end{theorem}

% Generator: Claude Opus 5 (High)
\begin{ainote}
Corsin states the chain rule without proof, as does the lecture at this point. The one-line
reason it is true is worth carrying, though: differentiability says $f$ and $g$ are each well
approximated by a linear map near the relevant point, and composing two approximations composes
the linear maps, with the error terms of the two compositions still $o(|h|)$. Making that
precise is the whole proof, and it is the same argument as in one variable.
\end{ainote}

% Supplement: Sascha Brack/Class Notes/Week_04_Notes_Friday_Updated.pdf, pp. 7--10
\begin{remark}[The chain rule as a sum over paths]
\label{rem:chain_rule_paths}
Stated as $D(g\circ f) = Dg\cdot Df$ the rule is compact but, as Sascha Brack puts it in his
notes, \qt{pretty abstract}. It says nothing about \emph{why} a sum appears when the rule is
written out in coordinates. His device is to draw the variables as a graph and read the formula
off it.

Write $x \in \mathbb{R}^n \xmapsto{\ f\ } y \in \mathbb{R}^m \xmapsto{\ g\ } z \in \mathbb{R}^k$
and join each variable to those it feeds into. Then $x_i$ influences $z_j$ only by first moving
some intermediate $y_\ell$, and the contribution of that route is the product of the two rates
along it. Summing over the routes gives the chain rule:

% Sonnet 5 (Medium) --- FIG-W04-03, after the variable graph on Sascha Brack's p. 10.
\begin{center}
\begin{tikzpicture}[>=Latex, scale=1.0]
  \node (xi) at (0,2.4) {$x_i$};
  \node (y1) at (-2.2,1.1) {$y_1$};
  \node (y2) at (-0.7,1.1) {$y_2$};
  \node (yd) at (0.75,1.1) {$\cdots$};
  \node (ym) at (2.2,1.1) {$y_m$};
  \node (zj) at (0,-0.2) {$z_j$};

  \foreach \t in {y1,y2,ym}{\draw[->, blue!70!black] (xi) -- (\t);}
  \foreach \s in {y1,y2,ym}{\draw[->, orange!90!black] (\s) -- (zj);}

  \node[blue!70!black, font=\scriptsize, above left] at (-1.25,1.75)
    {$\dfrac{\partial y_\ell}{\partial x_i}$};
  \node[orange!90!black, font=\scriptsize, below left] at (-1.25,0.55)
    {$\dfrac{\partial z_j}{\partial y_\ell}$};
\end{tikzpicture}
\end{center}

\[\frac{\partial z_j}{\partial x_i}
  \;=\; \sum_{\ell=1}^{m}
  \underbrace{\frac{\partial z_j}{\partial y_\ell}}_{(\Jac g)_{j\ell}}
  \underbrace{\frac{\partial y_\ell}{\partial x_i}}_{(\Jac f)_{\ell i}}
  \;=\; \big(\Jac g(f(x))\,\Jac f(x)\big)_{ji}.\]

The right-hand equality is just the definition of matrix multiplication, and that is precisely the
point: \textbf{matrix multiplication is a sum over paths}, and the chain rule looks like a product of
matrices precisely because composing functions composes those graphs. Sascha's summary is that
computing $\partial z_j/\partial x_i$ this way is nothing more than \qt{computing single
components of the chain rule}.

The reading is worth having because it makes the coordinate form reconstructible instead of
memorised: count the intermediate variables, one term each, each term a product of two factors,
and the index that is summed over is the one belonging to the middle layer. It is also the
form in which the rule is actually used in physics. Compare
\Cref{sec:chain_rule_along_path}, which is this picture with a single intermediate layer
$\gamma_1,\dots,\gamma_n$ and a single source variable $t$.
\end{remark}

% Generator: Claude Opus 5 (High) --- not in any source read so far; added because it is the
% cleanest one-line payoff of the chain rule along a path and is quoted constantly in physics.
\begin{theorem}[Euler's homogeneous function theorem]
\label{thm:eulers_homogeneous_function_theorem}
Let $f \in C^1(\mathbb{R}^n\setminus\{0\},\mathbb{R})$ be \newterm{positively homogeneous of
degree $k$} for some $k \in \mathbb{R}$, meaning that
\[f(tx) = t^kf(x) \qquad \text{for all } x \in \mathbb{R}^n\setminus\{0\} \text{ and all } t>0.\]
Then $\langle x, \nabla f(x)\rangle = k\,f(x)$ for all $x \in \mathbb{R}^n\setminus\{0\}$.
\end{theorem}

\begin{proof}
Fix $x \neq 0$ and read both sides of the homogeneity identity as functions of $t>0$ alone. The
left-hand side is $f$ composed with the path $\gamma(t) := tx$, whose velocity is
$\gamma'(t) = x$, so the chain rule along a path (\cref{sec:chain_rule_along_path}) gives
\[\frac{d}{dt}f(tx) = \big\langle \nabla f(tx),\ x\big\rangle.\]
On the right-hand side $f(x)$ is a constant in $t$, and therefore
$\frac{d}{dt}\big(t^kf(x)\big) = kt^{k-1}f(x)$. The two derivatives agree for every $t>0$;
evaluating at $t=1$ leaves $\langle \nabla f(x), x\rangle = k f(x)$, which is the claim by
symmetry of the inner product.
\end{proof}

\begin{aiexample}[Chain Rule for Vector Fields]
% Generator: Gemini 3.6 Flash (Medium)
Let $f(u,v) := u^2 + v^2$ and let $g(r,\theta) := (r\cos\theta, r\sin\theta)\transp$. The composite function is $(f \circ g)(r,\theta) = r^2\cos^2\theta + r^2\sin^2\theta = r^2$.
By the multi-variable Chain Rule, the Jacobian matrix of $f \circ g$ satisfies:
\[
D(f \circ g)(r,\theta) = Df(g(r,\theta)) \cdot Dg(r,\theta) = \begin{pmatrix} 2r\cos\theta & 2r\sin\theta \end{pmatrix} \begin{pmatrix} \cos\theta & -r\sin\theta \\ \sin\theta & r\cos\theta \end{pmatrix} = \begin{pmatrix} 2r & 0 \end{pmatrix},
\]
which matches the direct calculation $\frac{\partial}{\partial r}(r^2) = 2r$ and $\frac{\partial}{\partial \theta}(r^2) = 0$.
\end{aiexample}

% Source: Corsin Nick/Class Notes/Week 4.pdf, pp. 5--6
\begin{exercise}[Jacobian of a composite map]
\label{ex:chain_rule_computation}
Compute $\Jac(g\circ f)$ for
\[f : \mathbb{R}^3\to\mathbb{R}^2,\ x \mapsto (x_1x_2,\ x_2+x_3), \qquad g : \mathbb{R}^2\to\mathbb{R}^2,\ x \mapsto (e^{x_1},\ x_1x_2).\]
\end{exercise}

% Kept inline: solved directly in Corsin Nick/Class Notes/Week 4.pdf, pp. 5--6.
\begin{exercisesolution}[\cref{ex:chain_rule_computation}]
We have
\[\Jac f(x) = \begin{pmatrix} x_2 & x_1 & 0 \\ 0 & 1 & 1\end{pmatrix}, \qquad \Jac g(x) = \begin{pmatrix} e^{x_1} & 0 \\ x_2 & x_1 \end{pmatrix}.\]
So:
\[
\begin{aligned}
\Jac(g\circ f) &= \begin{pmatrix} e^{x_1x_2} & 0 \\ x_2+x_3 & x_1x_2\end{pmatrix}\begin{pmatrix} x_2 & x_1 & 0 \\ 0 & 1 & 1\end{pmatrix} \\
&= \begin{pmatrix} x_2e^{x_1x_2} & x_1e^{x_1x_2} & 0 \\ x_2^2 + x_2x_3 & 2x_1x_2 + x_1x_3 & x_1x_2 \end{pmatrix}
\end{aligned}
\]
\end{exercisesolution}

% Source: Jérôme Paschoud/Notizen/Woche 4.2 Kettenregel.pdf, p. 1
\begin{exercise}[Computing differentials with and without the chain rule]
\label{ex:chain_rule_computation_paschoud}
Let $f: \mathbb{R}^2 \to \mathbb{R}^3$ and $g: \mathbb{R}^3 \to \mathbb{R}$ be defined by
\[f(x,y) := \begin{pmatrix} x+y \\ x-y \\ x^2+y^2-1 \end{pmatrix}, \qquad g(x,y,z) := x^2+y^2+z^2.\]
Compute the differential $D(g \circ f)$ in two ways:
\begin{enumerate}[label=\textbf{(\alph*)}]
  \item By directly computing the composition $g \circ f$ and then differentiating.
  \item By computing $Df$ and $Dg$ and using the chain rule.
\end{enumerate}
\end{exercise}

\begin{exercisesolution}[\cref{ex:chain_rule_computation_paschoud}]
\begin{enumerate}[label=\textbf{(\alph*)}]
  \item \textbf{Directly:} First compute $g(f(x,y))$:
    \begin{align*}
    (g \circ f)(x,y) &= (x+y)^2 + (x-y)^2 + (x^2+y^2-1)^2 \\
    &= (x^2+2xy+y^2) + (x^2-2xy+y^2) + (x^4+y^4+1+2x^2y^2-2x^2-2y^2) \\
    &= 2x^2 + 2y^2 + x^4 + y^4 + 1 + 2x^2y^2 - 2x^2 - 2y^2 \\
    &= x^4 + y^4 + 2x^2y^2 + 1 \\
    &= (x^2+y^2)^2 + 1.
    \end{align*}
    Now differentiate this scalar function with respect to $x$ and $y$:
    \[D(g \circ f)(x,y) = \begin{pmatrix} 2(x^2+y^2)\cdot 2x & 2(x^2+y^2)\cdot 2y \end{pmatrix} = \begin{pmatrix} 4x(x^2+y^2) & 4y(x^2+y^2) \end{pmatrix}.\]

  \item \textbf{Via chain rule:} Compute the differentials of $f$ and $g$ individually.
    \[Df(x,y) = \begin{pmatrix} 1 & 1 \\ 1 & -1 \\ 2x & 2y \end{pmatrix}, \qquad Dg(u,v,w) = \begin{pmatrix} 2u & 2v & 2w \end{pmatrix}.\]
    Evaluate $Dg$ at $f(x,y)$:
    \[Dg(f(x,y)) = \begin{pmatrix} 2(x+y) & 2(x-y) & 2(x^2+y^2-1) \end{pmatrix}.\]
    Multiply the matrices:
    \begin{align*}
    D(g \circ f)(x,y) &= Dg(f(x,y)) Df(x,y) \\
    &= \begin{pmatrix} 2(x+y) & 2(x-y) & 2(x^2+y^2-1) \end{pmatrix} \begin{pmatrix} 1 & 1 \\ 1 & -1 \\ 2x & 2y \end{pmatrix} \\
    &= \begin{pmatrix} 4x + 4x(x^2+y^2-1) & 4y + 4y(x^2+y^2-1) \end{pmatrix} \\
    &= \begin{pmatrix} 4x(x^2+y^2) & 4y(x^2+y^2) \end{pmatrix}.
    \end{align*}
    Both methods yield the same result.
\end{enumerate}
\end{exercisesolution}

% Source: old_exams/examHS24.pdf (13 February 2025), Wahr/Falsch-Fragen, Aufgabe 3, p. 2
% Extractor: Gemini 3.1 Pro (High)
\begin{exercise}[True or False \textnormal{(the higher-order chain rule)}]
\label{ex:examHS24_TF_chain_rule}
Let $u \in C^2(\mathbb{R}^n)$ and $F = (F_1, \dots, F_n) \in C^2(\mathbb{R}^n, \mathbb{R}^n)$, with
$n \ge 2$. Decide whether each of the following identities is true or false.
\begin{enumerate}[label=\textbf{(\alph*)}]
  \item The identity $\partial_i(u \circ F) = \sum_{j=1}^n \partial_ju \circ F \, \partial_iF_j$ holds for all $i \in \{1, \dots, n\}$.
  \item The identity $\partial_{ii}(u \circ F) = \sum_{j=1}^n \partial_{jj}u \circ F \, (\partial_iF_j)^2 + \sum_{j=1}^n \partial_ju \circ F \, \partial_{ii}F_j$ holds for all $i \in \{1, \dots, n\}$.
  \item The identity $\partial_i\partial_j(u \circ F) = \partial_j\partial_i(u \circ F)$ holds for all $i, j \in \{1, \dots, n\}$.
\end{enumerate}
\exinfo{This exercise is taken from the exam of 13 February 2025 (Prof.\ Joaquim Serra), where it appears as the third
block of true/false questions.}
\exsol{sol:chain_rule_higher_order_tf}
\end{exercise}

% Source: old_exams/examFS24.pdf (21 August 2024), Wahr/Falsch-Fragen, Aufgabe 3, pp. 2--3
% Extractor: Claude Opus 5 (High)
% Presented as a worked example rather than as an exercise: the block is three one-line verdicts,
% and its value is entirely in seeing which hypothesis each one consumes. Its companion above,
% ex:examHS24_TF_chain_rule, is the same theme set as a problem, and assumes F in C^2 where this
% one assumes only C^1 --- which is the whole of part (c).
\begin{example}[Which identities are forced, and by what]
\label{ex:forced_identities_chain_rule}
% Generator: Claude Opus 5 (High) --- the three statements are the examiner's; the discussion is
% written for these notes.
Let $u \in C^2(\mathbb{R}^n)$ and let $F = (F_1,\dots,F_n) \in C^1(\mathbb{R}^n,\mathbb{R}^n)$,
with $n \geq 2$. Exactly one of the following three identities is forced by those hypotheses, and
the two that fail do so for quite different reasons.
\begin{enumerate}[label=\textbf{(\alph*)}]
  \item \textbf{$\partial_iF_j = \partial_jF_i$ for all $i,j$: not forced.} This asks the Jacobi
    matrix $\Jac F$ to be symmetric at every point, which is the integrability condition of
    \cref{lem:integrability_conditions}: it holds precisely when $F$ is a candidate for a gradient,
    and nothing in the hypotheses makes it one. The rotation field
    $F(x) := (-x_2,\, x_1,\, 0,\dots,0)$ is as smooth as one likes and has $\partial_1F_2 = 1$
    while $\partial_2F_1 = -1$.
  \item \textbf{$\partial_j\partial_iu = \partial_i\partial_ju$ for all $i,j$: forced.} This is
    Schwarz's lemma (\cref{thm:schwarz_mixed_partials}), and $u \in C^2$ is exactly its hypothesis.
    Note how differently the two statements read once they are set side by side: \textbf{(a)}
    constrains an arbitrary vector field, \textbf{(b)} constrains the gradient of a function. Only
    the second is automatic.
  \item \textbf{$u \circ F \in C^2$: not forced.} The chain rule gives
    $\partial_i(u\circ F) = \sum_{j}(\partial_ju\circ F)\,\partial_iF_j$, and the factors
    $\partial_iF_j$ are merely continuous, so this says only that $u \circ F \in C^1$. To
    differentiate once more one would have to differentiate $\partial_iF_j$, which the hypotheses
    do not permit. Take $n = 2$, $u(y_1,y_2) := y_1$ and $F(x_1,x_2) := (x_1\lvert x_1\rvert,\, x_2)$.
    Then $F \in C^1$, since $\partial_1F_1 = 2\lvert x_1\rvert$ is continuous, and
    $(u\circ F)(x) = x_1\lvert x_1\rvert$ has the same non-differentiable first derivative at
    $x_1 = 0$. Consequently $u \circ F \notin C^2$.
\end{enumerate}
The general rule that \textbf{(c)} illustrates is that a composition inherits the \emph{weaker} of
the two regularities. Raise $F$ to $C^2$ and $u \circ F$ becomes $C^2$, which is what
\cref{ex:examHS24_TF_chain_rule}\textbf{(c)} is then free to assume.
\exinfo{This example is taken from the exam of 21 August 2024 (Prof.\ Joaquim Serra), where it appears as the third block
of true/false questions.}
\end{example}
